% NE 630 -- Lesson 21 handout
% Compile: latexmk -xelatex -interaction=nonstopmode -halt-on-error lesson_21_handout.tex
% This source also compiles with pdflatex when using the built-in fallback.
%
% Use the course's shared reading-guide interface when it is available.
% The embedded fallback keeps this source portable outside the repository.
% To force the fallback, replace the next line with \def\NEOriginalStyle{0}.
\IfFileExists{ne630_ch04_reading_guide_common.tex}{\def\NEOriginalStyle{1}}{\def\NEOriginalStyle{0}}
\ifnum\NEOriginalStyle=1\relax
  % NE 630 -- shared reading-guide definitions, Lessons 15--21
% Included by lesson_15_handout.tex through lesson_21_handout.tex.
% Layout/palette adapted from the user's ne630boardhandout class.
% Content: supplied chapter(1).tex and lecture_15.tex through lecture_19.tex.
% The supplied chapter numbers the four-factor discussion as Section 4.5.
% Drawings below adapt supplied Fig. 4.1, Fig. 4.5, Fig. 4.6, and the L19
% pin-cell drawing. They are qualitative schematics, not simulation results.
\documentclass[10pt,letterpaper]{article}
\usepackage{amsmath,mathtools}
\usepackage{libertinus,microtype}
\usepackage[left=0.48in,right=0.48in,top=0.48in,bottom=0.44in,
  headheight=20pt,headsep=5pt,footskip=17pt,includeheadfoot]{geometry}
\usepackage{graphicx,xcolor,tabularx,booktabs,array,ragged2e}
\usepackage{tikz}
\usetikzlibrary{arrows.meta,calc,positioning,patterns}
\usepackage[most]{tcolorbox}
\usepackage{fancyhdr,lastpage}
\usepackage[hidelinks]{hyperref}
\hypersetup{pdfauthor={NE 630},pdfsubject={The Power Reactor Core}}
\definecolor{KSUPurple}{HTML}{512888}
\definecolor{KSUDeepPurple}{HTML}{3B1B64}
\definecolor{KSUOrange}{HTML}{CA7C1B}
\definecolor{KSUGray}{HTML}{5D5D5D}
\setlength{\parindent}{0pt}
\setlength{\parskip}{3pt}
\setlength{\emergencystretch}{1.5em}
\setlength{\abovedisplayskip}{4pt}
\setlength{\belowdisplayskip}{4pt}
\setlength{\abovedisplayshortskip}{2pt}
\setlength{\belowdisplayshortskip}{3pt}
\setlength{\jot}{3pt}
\renewcommand{\arraystretch}{1.12}
\newcolumntype{Y}{>{\RaggedRight\arraybackslash}X}
\newcommand{\LessonNo}{15}
\pagestyle{fancy}
\fancyhf{}
\fancyhead[L]{\sffamily\bfseries\small NE 630}
\fancyhead[C]{\sffamily\small The Power Reactor Core}
\fancyhead[R]{\sffamily\bfseries\small Lesson \LessonNo}
\fancyfoot[L]{\sffamily\scriptsize  }
\fancyfoot[C]{\sffamily\scriptsize \thepage\ of \pageref*{LastPage}}
\fancyfoot[R]{\sffamily\scriptsize }
\renewcommand{\headrulewidth}{0.5pt}
\makeatletter
\renewcommand{\headrule}{\hbox to\headwidth{\color{KSUPurple}\leaders\hrule height .5pt\hfill}}
\makeatother
\newcommand{\Lesson}[4]{%
  \renewcommand{\LessonNo}{#1}%
  \noindent\begin{minipage}[c]{0.025\linewidth}\color{KSUPurple}\rule{3.5pt}{32pt}\end{minipage}%
  \begin{minipage}[c]{0.815\linewidth}
   {\sffamily\bfseries\fontsize{18}{20}\selectfont #2}\par
   {\sffamily\fontsize{9}{10.5}\selectfont #4}
  \end{minipage}\hfill
  \begin{minipage}[c]{0.14\linewidth}\raggedleft
    {\sffamily\color{KSUPurple}\bfseries\fontsize{25}{26}\selectfont #1}\par
    {\sffamily\scriptsize LESSON}
  \end{minipage}\par\vspace{3pt}
  \begin{tcolorbox}[enhanced,boxrule=.45pt,arc=1pt,colback=KSUPurple!4,
   colframe=KSUPurple!40,left=5pt,right=5pt,top=3pt,bottom=3pt,
   before skip=0pt,after skip=6pt]
   \sffamily\fontsize{9}{10.5}\selectfont\textbf{Read:} #3
  \end{tcolorbox}}
\newcommand{\Sec}[2]{\par\addvspace{5pt}\noindent
  \colorbox{KSUPurple}{\color{white}\sffamily\bfseries\fontsize{9}{10}\selectfont\strut #1}
  \hspace{3pt}{\sffamily\bfseries\color{KSUDeepPurple}\fontsize{10.2}{11.5}\selectfont #2}\par\vspace{1pt}}
\newcommand{\eqread}[1]{{\sffamily\fontsize{7.9}{9}\selectfont\color{KSUGray}#1}}
\newcommand{\figcap}[1]{\par\vspace{1pt}{\sffamily\fontsize{7.7}{9}\selectfont\color{KSUGray}#1}\par}
\newtcolorbox{KeyBox}[1]{enhanced,boxrule=.45pt,arc=1pt,
 colback=KSUPurple!3,colframe=KSUPurple!40,left=5pt,right=5pt,top=4pt,bottom=4pt,
 before skip=4pt,after skip=4pt,before upper=\RaggedRight,fonttitle=\sffamily\bfseries\small,
 coltitle=KSUDeepPurple,title={#1},colbacktitle=KSUPurple!7,
 titlerule=0pt,toptitle=2pt,bottomtitle=2pt}
\newtcolorbox{Explore}[1]{enhanced,boxrule=.5pt,arc=1pt,
 colback=KSUOrange!5,colframe=KSUOrange!65,left=6pt,right=6pt,top=4pt,bottom=4pt,
 before skip=7pt,after skip=3pt,before upper=\RaggedRight,
 title={OpenMC \enspace\textbar\enspace #1},fonttitle=\sffamily\bfseries\fontsize{9}{10}\selectfont,
 coltitle=KSUDeepPurple,colbacktitle=KSUOrange!10,titlerule=0pt,toptitle=2pt,bottomtitle=2pt}
\newcommand{\SourceLine}[1]{\par\vspace{2pt}{\sffamily\fontsize{7.2}{8.4}\selectfont\color{KSUGray}#1}\par}
\newcommand{\colstart}{\begin{minipage}[t]{.482\linewidth}\vspace{0pt}\RaggedRight}
\newcommand{\colnext}{\end{minipage}\hfill\begin{minipage}[t]{.482\linewidth}\vspace{0pt}\RaggedRight}
\newcommand{\colend}{\end{minipage}\par}
\newcommand{\uunit}[1]{\,\mathrm{#1}}
\newcommand{\metal}{\mathrm{metal}}
\newcommand{\water}{\mathrm{w}}
\newcommand{\fiiso}{\mathrm{fi}}
\newcommand{\feiso}{\mathrm{fe}}
\newcommand{\nucl}[2]{\ensuremath{{}^{#1}\mathrm{#2}}}



% Five reactor motifs, adapted from the supplied Fig. 4.1, with a compact TRISO zoom.
\newcommand{\FamilyFigure}{%
\begin{tikzpicture}[x=1cm,y=1cm,line width=.42pt,font=\sffamily\fontsize{7.5}{8.5}\selectfont,
 every node/.style={inner sep=1pt},>=Latex]
\begin{scope}[shift={(0,0)}]
 \draw[fill=black!3] (-.60,-.50) rectangle (.60,.50);
 \draw[fill=black!35] (0,0) circle (.34);
 \draw[fill=KSUPurple!27] (0,0) circle (.29);
 \node at (0,0) {Fuel}; \node at (0,.42) {H$_2$O};
 \node[align=center] at (0,-.84) {\textbf{LWR}\\square pin lattice};
\end{scope}
\begin{scope}[shift={(3.05,0)}]
 \draw[fill=black!3] (90:.67)--(30:.67)--(-30:.67)--(-90:.67)--(-150:.67)--(150:.67)--cycle;
 \draw[fill=black!35] (0,0) circle (.38);
 \draw[fill=KSUPurple!27] (0,0) circle (.32);
 \node at (0,0) {Fuel}; \node at (0,.50) {Na};
 \node[align=center] at (0,-.84) {\textbf{SFR}\\triangular pin lattice};
\end{scope}
\begin{scope}[shift={(6.10,0)}]
 \draw[fill=black!7] (-.78,-.50) rectangle (.78,.50);
 \draw[fill=white] (0,0) circle (.39);
 \draw (0,0) circle (.34);
 \foreach \ang in {0,60,...,300}{\draw[fill=KSUPurple!27] (\ang:.21) circle (.07);} 
 \draw[fill=KSUPurple!27] (0,0) circle (.07);
 \node[align=center] at (-.58,0) {D$_2$O};
 \node[align=center] at (0,-.84) {\textbf{PHWR / CANDU}\\bundle in a pressure tube};
\end{scope}
\begin{scope}[shift={(9.25,0)}]
 \draw[fill=black!12] (-.84,-.50) rectangle (.84,.50);
 \foreach \x/\y in {-.50/.25,0/.25,.50/.25,-.25/-.20,.25/-.20}{\draw[fill=white] (\x,\y) circle (.10);} 
 \foreach \x/\y in {-.56/-.18,-.19/.03,.19/.03,.56/-.18}{
   \draw[fill=KSUPurple!10,pattern=dots,pattern color=KSUPurple!70] (\x,\y) circle (.085);
   \draw (\x,\y) circle (.085);
 }
 \node[anchor=west] at (.86,.22) {He};
 \draw[->] (.82,.22)--(.59,.22);
 \node[anchor=west] at (.86,-.18) {TRISO};
 \draw[->] (.82,-.18)--(.62,-.18);
 \node[align=center] at (0,-.84) {\textbf{HTGR}\\TRISO in graphite};
\end{scope}
\begin{scope}[shift={(12.55,0)}]
 \draw[fill=KSUPurple!70] (0,0) circle (.10);  % kernel
 \draw[fill=KSUPurple!18] (0,0) circle (.18);  % buffer
 \draw[fill=gray!25] (0,0) circle (.24);       % inner PyC
 \draw[fill=green!12!black!8] (0,0) circle (.30); % SiC
 \draw[fill=gray!8] (0,0) circle (.36);        % outer PyC
 \foreach \r in {.10,.18,.24,.30,.36}{\draw (0,0) circle (\r);} 
 \node[font=\tiny] at (0,0) {K};
 \node[font=\tiny,anchor=west] at (.62,.25) {OPyC}; \draw[->] (.58,.24)--(.33,.24);
 \node[font=\tiny,anchor=west] at (.62,.06) {SiC};  \draw[->] (.58,.05)--(.28,.05);
 \node[font=\tiny,anchor=west] at (.62,-.12) {IPyC};\draw[->] (.58,-.11)--(.21,-.11);
 \node[font=\tiny,anchor=west] at (.62,-.30) {buffer};\draw[->] (.58,-.29)--(.14,-.20);
 \node[font=\tiny,anchor=east] at (-.62,.32) {kernel};\draw[->] (-.58,.30)--(-.08,.08);
 \node[align=center] at (0,-.84) {\textbf{TRISO particle}\\coated-fuel cross section};
\end{scope}
\end{tikzpicture}}

% Square and hexagonal pin cells, adapted from the L19 cell figure and Fig. 4.1.
\newcommand{\PinCellsFigure}{%
\begin{tikzpicture}[x=1cm,y=1cm,font=\sffamily\fontsize{8.5}{10}\selectfont,
 line width=.55pt,>=Latex,every node/.style={inner sep=1pt}]
\begin{scope}[shift={(0,0)}]
 \draw[fill=black!3] (-.88,-.88) rectangle (.88,.88);
 \draw[fill=KSUOrange!35] (0,0) circle (.58);
 \draw[fill=KSUPurple!25] (0,0) circle (.49);
 \draw[->] (0,0)--(45:.49) node[midway,above left=-1pt] {$r_f$};
 \draw[->] (0,0)--(-35:.58);\node[anchor=west] at (.57,-.40) {$r_c$};
 \draw[|<->|] (-.88,-1.10)--(.88,-1.10) node[midway,below=3pt] {$P$};
 \node at (0,1.10) {Square};
\end{scope}
\begin{scope}[shift={(3.80,0)}]
 \draw[fill=black!3] (90:1.016)--(30:1.016)--(-30:1.016)--(-90:1.016)--(-150:1.016)--(150:1.016)--cycle;
 \draw[fill=KSUOrange!35] (0,0) circle (.58);
 \draw[fill=KSUPurple!25] (0,0) circle (.49);
 \draw[densely dashed] (0,0)--(.88,0);
 \node[above] at (.69,.05) {$P/2$};
 \draw[|<->|] (-.88,-1.22)--(.88,-1.22) node[midway,below=3pt] {$P$};
 \node at (0,1.24) {Hexagonal};
 \node[anchor=west] at (.82,.77) {$a$};
\end{scope}
\end{tikzpicture}}

% Bookkeeping cycle adapted from supplied Fig. 4.5 (not a spatial plot).
\newcommand{\CycleFigure}{%
\begin{tikzpicture}[x=1cm,y=1cm,>=Latex,
 box/.style={draw=KSUPurple!45,fill=KSUPurple!3,rounded corners=1pt,
 minimum width=3.0cm,minimum height=1.05cm,align=center,inner sep=3pt},
 every node/.style={font=\sffamily\fontsize{8.2}{9.5}\selectfont}]
 \node[box] (n) at (0,0) {$n$\\thermal-fission births};
 \node[box] (en) at (3.85,0) {$\epsilon n$\\all fission births};
 \node[box] (epn) at (7.7,0) {$\epsilon p n$\\reach thermal energies};
 \node[box] (epfn) at (11.55,0) {$\epsilon p f n$\\absorbed in fuel};
 \node[box] (kn) at (15.4,0) {$\epsilon p f\eta_T n=k_\infty n$\\next thermal-fission births};
 \draw[->,line width=.6pt] (n)--(en) node[midway,above=2pt] {$\epsilon$};
 \draw[->,line width=.6pt] (en)--(epn) node[midway,above=2pt] {$p$};
 \draw[->,line width=.6pt] (epn)--(epfn) node[midway,above=2pt] {$f$};
 \draw[->,line width=.6pt] (epfn)--(kn) node[midway,above=2pt] {$\eta_T$};
\end{tikzpicture}}

% Qualitative Doppler/space-energy shielding sketch, adapted from Fig. 4.6.
\newcommand{\ResonanceFigure}{%
\begin{tikzpicture}[x=1cm,y=1cm,>=Latex,
 ax/.style={->,line width=.5pt,draw=black!75},
 cold/.style={line width=.85pt,draw=KSUPurple},
 hot/.style={line width=.85pt,draw=KSUOrange,dashed},
 every node/.style={font=\sffamily\fontsize{8.1}{9}\selectfont,inner sep=1pt}]
\begin{scope}
 \draw[ax] (0,0)--(3.5,0) node[right] {$E$};
 \draw[ax] (0,0)--(0,1.9) node[above] {$\sigma_a(E)$};
 \draw[cold] plot[smooth] coordinates {(0,.01)(.5,.04)(.9,.15)(1.2,.52)(1.4,1.2)(1.55,1.7)(1.7,1.2)(1.95,.43)(2.3,.10)(2.8,.03)(3.2,.01)};
 \draw[hot] plot[smooth] coordinates {(0,.02)(.5,.10)(.9,.28)(1.2,.72)(1.55,1.08)(1.9,.75)(2.3,.30)(2.8,.10)(3.2,.03)};
 \node[anchor=west,text=KSUPurple] at (.08,1.63) {$T_1$};
 \node[anchor=west,text=KSUOrange] at (2.26,1.09) {$T_2>T_1$};
\end{scope}
\begin{scope}[shift={(4.25,0)}]
 \draw[ax] (0,0)--(3.5,0) node[right] {$E$};
 \draw[ax] (0,0)--(0,1.9) node[above] {$\varphi_f/\varphi_m$};
 \draw[black!40,densely dotted] (0,1.5)--(3.2,1.5);
 \node[anchor=east] at (-.05,1.5) {$1$};
 \draw[cold] plot[smooth] coordinates {(0,1.5)(.6,1.49)(1,1.3)(1.25,.65)(1.5,.3)(1.75,.59)(2,1.26)(2.5,1.49)(3.2,1.5)};
 \draw[hot] plot[smooth] coordinates {(0,1.5)(.4,1.43)(.8,1.22)(1.2,.88)(1.55,.65)(1.9,.83)(2.3,1.18)(2.8,1.43)(3.2,1.49)};
\end{scope}
\end{tikzpicture}}

\else
  \documentclass[10pt,letterpaper]{article}
  \usepackage[margin=0.50in,top=0.47in,bottom=0.48in,footskip=0.23in]{geometry}
  \usepackage[T1]{fontenc}
  \usepackage{newtxtext,newtxmath}
  \usepackage{microtype}
  \usepackage{amsmath,booktabs,array,graphicx,xcolor}
  \usepackage{ragged2e}
  \usepackage{tikz}
  \usepackage[most]{tcolorbox}
  \usepackage{fancyhdr}
  \usepackage[hidelinks]{hyperref}
  \definecolor{ksupurple}{HTML}{512888}
  \definecolor{orange}{HTML}{CA7C1B}
  \setlength{\parindent}{0pt}
  \setlength{\parskip}{2pt}
  \setlength{\emergencystretch}{1em}
  \pagestyle{fancy}
  \fancyhf{}
  \renewcommand{\headrulewidth}{0pt}
  \renewcommand{\footrulewidth}{0.3pt}
  \fancyfoot[L]{\sffamily\fontsize{8}{9}\selectfont NE 630\enspace /\enspace Nuclear Reactor Theory}
  \fancyfoot[R]{\sffamily\fontsize{8}{9}\selectfont Lesson 21\enspace /\enspace \thepage}
  \newcommand{\Lesson}[4]{%
    {\sffamily\bfseries\fontsize{9}{11}\selectfont\color{ksupurple}NE 630\enspace |\enspace LESSON #1}\par
    \vspace{1pt}
    {\sffamily\bfseries\fontsize{19}{22}\selectfont #2\par}
    \vspace{2pt}
    {\fontsize{8.5}{10}\selectfont\textbf{Reading:} #3\par}
    {\fontsize{10}{12}\selectfont\color{ksupurple}\textit{#4}\par}
    \vspace{2pt}{\color{ksupurple}\hrule height 0.7pt}\vspace{5pt}}
  \newcommand{\Sec}[2]{\par\addvspace{6pt}%
    {\sffamily\bfseries\color{ksupurple}#1\enspace #2\par}\vspace{2pt}}
  \newcommand{\colstart}{\par\noindent\begin{minipage}[t]{0.483\textwidth}\vspace{0pt}\RaggedRight}
  \newcommand{\colnext}{\end{minipage}\hfill\begin{minipage}[t]{0.483\textwidth}\vspace{0pt}\RaggedRight}
  \newcommand{\colend}{\end{minipage}\par}
  \newtcolorbox{KeyBox}[1]{enhanced,boxrule=0.45pt,arc=1mm,
    colback=ksupurple!2,colframe=ksupurple!65,
    colbacktitle=ksupurple!9,coltitle=ksupurple,
    fonttitle=\sffamily\bfseries\fontsize{9}{10}\selectfont,title={#1},
    left=5pt,right=5pt,top=3pt,bottom=3pt,
    toptitle=3pt,bottomtitle=3pt,before skip=5pt,after skip=5pt}
  \newtcolorbox{Explore}[1]{enhanced,boxrule=0.45pt,arc=1mm,
    colback=orange!4,colframe=orange!75,
    colbacktitle=orange!12,coltitle=black,
    fonttitle=\sffamily\bfseries\fontsize{9}{10}\selectfont,title={#1},
    left=6pt,right=6pt,top=3pt,bottom=3pt,
    toptitle=3pt,bottomtitle=3pt,before skip=7pt,after skip=0pt}
  \newcommand{\figcap}[1]{{\fontsize{8}{9.5}\selectfont #1\par}}
\fi

% Extra, local helpers; the L18 interface above is otherwise unchanged.
\usepackage{amsmath,booktabs,array,graphicx,xcolor,tikz,ragged2e}
\usetikzlibrary{arrows.meta,calc,decorations.pathreplacing,positioning}
\providecolor{ksupurple}{HTML}{512888}
\providecolor{orange}{HTML}{CA7C1B}
\newcommand{\NEContinue}[2]{%
  {\sffamily\bfseries\fontsize{12}{14}\selectfont\color{ksupurple}Lesson #1\enspace |\enspace #2\par}
  \vspace{4pt}{\color{ksupurple}\hrule height 0.6pt}\vspace{5pt}}
\newcommand{\NEEx}[1]{\par\addvspace{4pt}\textbf{#1}\par}
\newcommand{\NECheck}{\textbf{Definition check.}\ }
\newcommand{\NEsmall}[1]{{\fontsize{8.3}{10}\selectfont #1\par}}
\newcommand{\NEbody}{%
  \fontsize{8.7}{10.25}\selectfont\RaggedRight
  \setlength{\parskip}{2pt}
  \setlength{\topsep}{2pt}\setlength{\partopsep}{0pt}
  \setlength{\abovedisplayskip}{4pt}\setlength{\belowdisplayskip}{4pt}
  \setlength{\abovedisplayshortskip}{3pt}\setlength{\belowdisplayshortskip}{3pt}}

% Editable reconstruction of the startup sketch. The source's rho_ex axis
% has been labelled as a budget to distinguish it from actual controlled rho.
\newcommand{\NEStartupFigure}{%
\begin{tikzpicture}[x=1cm,y=1cm,>=Stealth,line width=0.5pt,
  every node/.style={font=\fontsize{7.2}{8.2}\selectfont},
  dot/.style={circle,fill=ksupurple,inner sep=1.6pt}]
  \draw[->] (0,0)--(7.9,0);
  \draw[->] (0,0)--(0,3.32);
  \node[rotate=90] at (-0.44,1.60) {reactivity budget};
  \node[anchor=east] at (-0.05,0) {$0$};
  \coordinate (A) at (0.65,2.95);
  \coordinate (B) at (2.25,2.27);
  \coordinate (C) at (4.35,1.52);
  \coordinate (D) at (6.60,0.72);
  \draw[dashed,ksupurple!70] (A)--(B)--(C)--(D);
  \foreach \P in {A,B,C,D}{\node[dot] at (\P) {};}
  \draw[decorate,decoration={brace,amplitude=3pt,mirror}]
    (2.47,2.27)--node[right=4pt,align=left]{shutdown margin\\$|\Delta\rho_{\rm sd}|$}(2.47,2.95);
  \draw[decorate,decoration={brace,amplitude=3pt,mirror}]
    (4.58,1.52)--node[right=4pt,align=left]{temperature defect\\$|\Delta\rho_{\rm td}|$}(4.58,2.27);
  \draw[decorate,decoration={brace,amplitude=3pt,mirror}]
    (6.82,0.72)--node[right=4pt,align=left]{power defect\\$|\Delta\rho_{\rm pd}|$}(6.82,1.52);
  \draw[decorate,decoration={brace,amplitude=3pt,mirror}]
    (7.50,0.04)--node[right=4pt,align=left]{variable\\control}(7.50,0.72);
  \foreach \xx in {0.65,2.25,4.35,6.60}{\draw (\xx,-.035)--(\xx,.035);}
  \node[below=4pt,align=center] at (.65,0) {cold\\shutdown};
  \node[below=4pt,align=center] at (2.25,0) {cold\\critical};
  \node[below=4pt,align=center] at (4.35,0) {hot zero\\power};
  \node[below=4pt,align=center] at (6.60,0) {hot full\\power};
  \draw[->] (2.1,-.83)--node[below]{startup sequence}(6.8,-.83);
\end{tikzpicture}}

% Editable schematic of the L21 Doppler-broadening sketch, not nuclear data.
% Separate vertical scales avoid conflating cross section and flux.
\newcommand{\NEDopplerFigure}{%
\begin{tikzpicture}[x=1cm,y=1cm,>=Stealth,line width=0.65pt,
  every node/.style={font=\fontsize{8}{9}\selectfont}]
  \draw[->] (0,2.25)--(0,4.18) node[above] {$\sigma_a(E)$};
  \draw[->] (0,2.25)--(6.7,2.25) node[right] {$E$};
  \draw[ksupurple,domain=.12:6.45,samples=140,smooth]
    plot (\x,{2.42+1.42*exp(-((\x-3.2)/.30)^2)});
  \draw[orange,dashed,domain=.12:6.45,samples=140,smooth]
    plot (\x,{2.42+.73*exp(-((\x-3.2)/.63)^2)});
  \draw[->] (0,0)--(0,1.70) node[above] {$\varphi_f(E)/\varphi_m(E)$};
  \draw[->] (0,0)--(6.7,0) node[right] {$E$};
  \draw[ksupurple,domain=.12:6.45,samples=140,smooth]
    plot (\x,{1.36-1.04*exp(-((\x-3.2)/.38)^2)});
  \draw[orange,dashed,domain=.12:6.45,samples=140,smooth]
    plot (\x,{1.36-.77*exp(-((\x-3.2)/.79)^2)});
  \draw[ksupurple] (.35,-.45)--(.95,-.45);
  \node[anchor=west] at (1.0,-.45) {cooler fuel};
  \draw[orange,dashed] (3.45,-.45)--(4.05,-.45);
  \node[anchor=west] at (4.1,-.45) {hotter fuel};
\end{tikzpicture}}

% PRIMARY SOURCE: L21.pdf, handwritten pages 1--5, including all three
% numbered examples, the product-rule derivation, FTC/MTC derivations,
% the Doppler sketch, the Monte Carlo caution, and the fast-reactor pointer.
% The source's near-critical -91 pcm/K result is retained, not replaced.
% An explicitly labelled definition check supplies the exact-reactivity
% finite-difference result. The inconsistent denominator in the handwritten
% intermediate step is reordered to match its numerator and negative slope.
% Some right-edge prose is physically cropped; uncertain words and the
% incomplete MTC range are not guessed.
\hypersetup{pdftitle={NE 630 -- Lesson 21: Reactivity coefficients},pdfauthor={Jeremy Roberts}}
\begin{document}
\NEbody
\Lesson{21}{Reactivity coefficients}
 {FNRP \S\,9.2--9.4; the handwritten derivations and examples for Lesson 21.}
 {How does changing one system parameter change reactivity?}
\textbf{Objective:} Quantify reactivity changes due to fuel temperature, moderator temperature, and other parameters, including coolant void fraction.

\colstart
\Sec{1}{From integral changes to local sensitivity}
Last time: $\rho=(k-1)/k$, its numerical units, and the integral reactivity changes from shutdown to full power. Now hold the other independent parameters fixed and define
$$
 \boxed{\alpha_x=\left.\frac{\partial\rho}{\partial x}
 \right|_{\rm all\ other\ parameters\ fixed}.}
$$
For $x=T_F$, this is the \textbf{fuel temperature coefficient}, $\alpha_F$ (FTC). For $x=T_M$, it is the \textbf{moderator temperature coefficient}, $\alpha_M$ (MTC). Temperature coefficients are commonly reported in pcm/K.  For boiling water reactors, the coolant void fraction (or steam quality) is another possible $x$.

\vspace{10pt}

Starting directly from the definition of $\rho$,
$$
 \frac{\partial\rho}{\partial k}=\frac1{k^2},\qquad
 \alpha_x=\frac1{k^2}\frac{\partial k}{\partial x}.
$$

Typically, coefficients are reported for near-critical states, or
$$\alpha_x\simeq\frac1k\frac{\partial k}{\partial x}\qquad(k\simeq1).$$
Away from $k=1$, the two normalizations differ.

\Sec{2}{A Finite Difference Approximation}
Suppose we modeled a PWR fuel pin at two different
moderator temperatures using OpenMC and found
the following multiplication factors:
$$
 \underbrace{k_{600}=1.28397}_{\rm nominal},\qquad
 \underbrace{k_{610}=1.27228}_{\rm perturbed}.
$$
To estimate $\alpha_M$ at the nominal $T_M=600$ K, we can apply a {\bf finite-difference} approximation, or
$$
 \alpha_M \approx \left.\frac{\partial k}{\partial T_M}\right|_{600}
 \simeq\frac{k_{610}-k_{600}}{610-600}
 =-0.001169\ \mathrm{K}^{-1} ,
$$
or, in the more common units,
$$
 \boxed{\alpha_M\simeq\frac{-0.001169}{1.28397}
 =-9.1046\times10^{-4}\ \mathrm{K}^{-1}
 \simeq-91\ \mathrm{pcm/K}.}
$$

\vspace{10pt}

\begin{KeyBox}{Assumption check}
Using $\rho=1-1/k$ at both temperatures instead gives
$$
 \frac{\rho_{610}-\rho_{600}}{10}
 =\frac{k_{610}-k_{600}}{10k_{600}k_{610}}
 =-71.6\ \mathrm{pcm/K}.
$$
Because $k\simeq1.28$, this differs materially from the near-critical result; the difference is not rounding!  Unless directed otherwise, the near-critical approximation should be used for reporting reactivity coefficients.
\end{KeyBox}

\colnext
\Sec{3}{Expressing $\alpha_x$ with the Four Factors}
Recall $k_\infty=\eta_Tfp\varepsilon$. The handwritten $\eta$ is written $\eta_T$ here, following the slide notation. The product rule gives
\begin{align*}
 \frac{\partial k_\infty}{\partial x}
 &=fp\varepsilon\frac{\partial\eta_T}{\partial x}
 +\eta_Tp\varepsilon\frac{\partial f}{\partial x}\\
 &\quad+\eta_Tf\varepsilon\frac{\partial p}{\partial x}
 +\eta_Tfp\frac{\partial\varepsilon}{\partial x}.
\end{align*}
Divide by $k_\infty=\eta_Tfp\varepsilon$:
\begin{align*}
 \frac1{k_\infty}\frac{\partial k_\infty}{\partial x}
 &=\frac1{\eta_T}\frac{\partial\eta_T}{\partial x}
 +\frac1f\frac{\partial f}{\partial x}\\
 &\quad+\frac1p\frac{\partial p}{\partial x}
 +\frac1\varepsilon\frac{\partial\varepsilon}{\partial x}.
\end{align*}
Thus, in the near-critical approximation used throughout the handwritten four-factor analysis,
$$
 \boxed{\alpha_x\simeq
 \frac1{\eta_T}\frac{\partial\eta_T}{\partial x}
 +\frac1f\frac{\partial f}{\partial x}
 +\frac1p\frac{\partial p}{\partial x}
 +\frac1\varepsilon\frac{\partial\varepsilon}{\partial x}.}
$$
Each term is the fractional sensitivity of one factor. The exact derivative of $\rho_\infty$ is $1/k_\infty$ times this sum.

\Sec{4}{Recall the UO$_2$ resonance model}
For a UO$_2$ fuel element,
$$
 p=\exp\!\left[-\frac{V_fN_{fe}I}{V_m\xi^m\bar\Sigma_s^m}\right].
$$
Here $V_f$ and $V_m$ are the fuel and moderator volumes, $N_{fe}$ is the fertile-nuclide number density in the fuel, and $\xi^m\bar\Sigma_s^m$ is the moderator slowing-down power. The effective heterogeneous resonance integral is $I$.

\vspace{10pt}

With $\rho_f$ the fuel mass density and $D$ the fuel diameter, define
\begin{align*}
 I(300)\,[\mathrm b]&=4.45+26.6\sqrt{\frac4{\rho_fD}},\\
 I(T_F)&=I(300)\big[1+\tilde\gamma
 (\sqrt{T_F}-\sqrt{300})\big],\\
 \tilde\gamma\,[\mathrm{K}^{-1/2}]&=0.0061+0.0047\frac4{\rho_fD}.
\end{align*}
Insert $\rho_f$ in g/cm$^3$, $D$ in cm, and temperatures in kelvin. The subscript on $\rho_f$ distinguishes mass density from reactivity.  These relations are equivalent to FNRP Eq.\ (4.40), Table 4.3, Eqs.\ (9.13)--(9.14), and Table 9.1.

\vspace{10pt}


At fixed moderator conditions, the explicit $T_F$ dependence in this simplified model enters through $I(T_F)$ and hence $p$. Moderator temperature affects moderator density, so it changes both $p$ and $f$. The next page differentiates those pathways separately.

\colend

\newpage
\NEContinue{21}{Fuel and moderator temperature feedback}
\colstart
\Sec{5}{Derive the fuel temperature coefficient}
With the moderator and geometry fixed, write
$$
 C_F=\frac{V_fN_{fe}}{V_m\xi^m\bar\Sigma_s^m},\qquad
 \ln p=-C_F I(T_F).
$$
Since $p$ is the only factor with explicit $T_F$ dependence in this model,
\begin{align*}
 \alpha_F&\simeq\frac1p\frac{\partial p}{\partial T_F}
 =-C_F\frac{\partial I}{\partial T_F}\\
 &=-\ln\!\left(\frac1p\right)\frac1I
 \frac{\partial I}{\partial T_F}.
\end{align*}
The empirical relation gives
$$
 \frac{\partial I}{\partial T_F}
 =I(300)\frac{\tilde\gamma}{2\sqrt{T_F}},\qquad
 C_F I(300)=\ln\!\left[\frac1{p(300)}\right].
$$
Therefore
$$
 \boxed{\alpha_F\simeq-\ln\!\left[\frac1{p(300)}\right]
 \frac{\tilde\gamma}{2\sqrt{T_F}}.}
$$
The negative sign follows from $0<p(300)<1$ and $\tilde\gamma>0$. Typical values for $\alpha_T$ are $\alpha_F\simeq-(2\text{--}3)$ pcm/K for PWRs and about $-10$ pcm/K for our TRIGA reactor.

\vspace{10pt}

Why does $I$ depend on $T_F$? {\bf Doppler broadening}!


% \begin{center}
% \resizebox{0.90\linewidth}{!}{\NEDopplerFigure}
% \end{center}
% \vspace{-5pt}
% \figcap{Reconstruction of the handwritten sketch: solid, cooler fuel; dashed, hotter fuel. Curves are schematic, with separate vertical scales; they are not evaluated cross-section data.}
% The broadened resonance and the associated flux depression must be considered together. As emphasized in the supplied L18 handout, absorption depends on the integral of $\sigma_a(E)\varphi_f(E)$, not just the peak cross section. Absorption in less-shielded resonance wings increases the effective integral $I$ and lowers $p$.
% % The prose to the right of the original L21 sketch is cut off. The complete
% % mechanism sentence above is attributed to the supplied L18 handout rather
% % than presented as a guessed transcription of that clipped annotation.

\begin{KeyBox}{Example: Fuel Temperature Coefficient}
For $p(300)=0.8$, compute $\alpha_F$ at $T_F=1200$ K for a UO$_2$ fuel
element with a diameter of 1 cm and mass density of 10 g/cm$^3$.
\begin{align*}
 \tilde\gamma
 &=0.0061+0.0047\frac4{(10)(1)}
 =0.00798\ \mathrm{K}^{-1/2},\\
 \alpha_F(1200)
 &\simeq-\ln(1/0.8)\frac{0.00798}{2\sqrt{1200}}\\
 &=-2.57\times10^{-5}\ \mathrm{K}^{-1}
 =-2.57\ \mathrm{pcm/K}\approx\boxed{-2.6\ \mathrm{pcm/K}}.
\end{align*}
\end{KeyBox}



\colnext
\Sec{6}{Derive the moderator temperature coefficient}
Recall
$$
 f=\frac1{1+\varsigma\displaystyle\frac{V_m}{V_f}
 \displaystyle\frac{\bar\Sigma^M_{a,2}}{\bar\Sigma^F_{a,2}}}.
$$
In the density-only model, $T_M$ changes the moderator number densities: decreasing $\rho_m=\rho_{\mathrm{H_2O}}$ lowers both $\bar\Sigma^M_{a,2}$ and $\xi^m\bar\Sigma_s^m$.

Define the volumetric thermal-expansion coefficient by
$$
 \beta_M=-\frac1{\rho_m}\frac{d\rho_m}{dT_M}.
$$
Careful: this $\beta_M$ is distinct from the delayed neutron fraction $\beta$.

\vspace{6pt}

\textbf{Thermal utilization.} Collect the fixed quantities into $C$ and write $f=[1+C\rho_m(T_M)]^{-1}$. Differentiating gives
\begin{align*}
 \frac1f\frac{df}{dT_M}
 &=-\frac{C\rho_m}{1+C\rho_m}
   \frac1{\rho_m}\frac{d\rho_m}{dT_M}\\
 &=(1-f)\beta_M>0.
\end{align*}

\vspace{6pt}

\textbf{Resonance escape.} Write $p=\exp[-A/\rho_m(T_M)]$, with $A$ fixed.  Then
\begin{align*}
 \frac1p\frac{dp}{dT_M}
 &=\frac{A}{\rho_m}\frac1{\rho_m}\frac{d\rho_m}{dT_M}\\
 &=-\beta_M\ln(1/p)<0.
\end{align*}
Combine these two contributions:
$$
 \boxed{\alpha_M\simeq-\beta_M\big[\ln(1/p)-(1-f)\big].}
$$
The \textbf{overall sign depends on their relative magnitudes}. In this model, negative moderator feedback requires $\ln(1/p)>(1-f)$.

\vspace{6pt}

\begin{KeyBox}{Example: Moderator Temperature Coefficient}
For $p=0.6$, $f=0.95$, and $\beta_M=2.9\times10^{-4}$ K$^{-1}$,
\begin{align*}
 \alpha_M&\simeq-2.9\times10^{-4}
 \big[\ln(1/0.6)-(1-0.95)\big]\\
 &=-1.3364\times10^{-4}\ \mathrm{K}^{-1}\\
 &=-13.36\ \mathrm{pcm/K}\approx\boxed{-13\ \mathrm{pcm/K}}.
\end{align*}
The resonance-escape penalty outweighs the utilization gain for these values.
\end{KeyBox}



\colend


\end{document}
